% Grado 8 -- generado por tools/build_tex.py a partir de raw/grado_08.txt
\section{Grado 8}\label{grado:8}

% PDF p. 59
\dba{1}{Reconoce la existencia de los números irracionales como números no racionales y los describe de acuerdo con sus características y propiedades.}

\evidencias

\begin{itemize}
  \item Utiliza procedimientos geométricos para representar números racionales e irracionales.
  \item Identifica las diferentes representaciones (decimales y no decimales) para argumentar por qué un número es o no racional.
\end{itemize}

\ejemplo

Marin a Juliá n Marcela “Yo em pecé Catalina deci puse el punt co n cero, pone m r al y em pe o de cons natu rales los nú meros cé a 0.12 34 ec utiva, de es forma 4…. y así 56789101 112131 decir, men su te cesiva-

En clase de matemáticas la profesora pidió a los estudiantes que construyeran números en su representación decimal. Algunos estudiantes hicieron lo siguiente:

Marina dijo: “Yo empecé con el cinco como entero. Para formar los decimales utilicé un dado y lo lancé 10 veces, y así formé el número con 10 dígitos después del punto”

Julián dijo: “Yo empecé con cero, puse el punto de decimal y empecé a poner los números naturales de forma consecutiva, es decir, 0.1234567891011121314… y así sucesivamente”

Catalina dijo: “Yo recordé algo del año pasado y lo formé dividiendo en la calculadora 1 entre 3”

Marcela dijo: “Yo me inspiré en lo que Marina hizo, pero mi número se formaría pensando en que siempre voy a poder seguir tirando el dado, por tanto tendría infinitos dígitos decimales”

Analiza la manera en que Marina, Julián, Catalina y Marcela construyeron sus números y argumenta cuáles de ellos serían racionales y cuáles no. Propone otras maneras de construir números y argumenta cuáles de ellos serían racionales y cuáles no.


% PDF p. 59
\dba{2}{Construye representaciones, argumentos y ejemplos de propiedades de los números racionales y no racionales.}

\evidencias

\begin{itemize}
  \item Utiliza procedimientos geométricos o aritméticos para construir algunos números irracionales y los ubica en la recta numérica.
  \item Justificar procedimientos con los cuales se representa geométricamente números racionales y números reales.
  \item Construye varias representaciones (geométrica, decimales o no decimales) de un mismo número racional o irracional.
\end{itemize}

\ejemplo

Si en la siguiente representación, el triángulo y el

\fig{g08_p59_2}{Expresiones con símbolos: √(△ + □) comparada con √△ + √□, donde △ y □ son números cuadrados.}

Asigna valores en las casillas y y utiliza la calculadora para establecer la existencia de números que hagan verdadera la igualdad. Argumenta este hecho y escribe una consecuencia que pueda inferirse a partir de esta exploración. Construye otras representaciones con productos, cocientes y potencias y analiza lo que sucede en cada caso.


% PDF p. 60
\dba{3}{Reconoce los diferentes usos y significados de las operaciones (convencionales y no convencionales) y del signo igual (relación de equivalencia e igualdad condicionada) y los utiliza para argumentar equivalencias entre expresiones algebraicas y resolver sistemas de ecuaciones.}

\evidencias

\begin{itemize}
  \item Reconoce el uso del signo igual como relación de equivalencia de expresiones algebraicas en los números reales.
  \item Propone y ejecuta procedimientos para resolver una ecuación lineal y sistemas de ecuaciones lineales y argumenta la validez o no de un procedimiento
  \item Usa el conjunto solución de una relación (de equivalencia y de orden) para argumentar la validez o no de un procedimiento.
\end{itemize}

\ejemplo

En clase de matemáticas el profesor pidió a los estudiantes analizar tres expresiones y hablar acerca de sus posibles relaciones. Las tres expresiones fueron: q x - 1 q 4x+1 q 5 Al respecto Carlos y José escribieron:

\fig{g08_p60_1}{Tarjetas con los argumentos algebraicos de Carlos y de José sobre una propiedad de los números.}

Analiza los escritos de Carlos y José y presenta argumentos que confirmen o refuten lo que ellos han hecho. Determina si José tiene razón al dudar si aplica o no la propiedad transitiva. ¿De qué depende que la pueda aplicar o no?


% PDF p. 60
\dba{4}{Describe atributos medibles de diferentes sólidos y explica relaciones entre ellos por medio del lenguaje algebraico.}

\evidencias

\begin{itemize}
  \item Utiliza lenguaje algebraico para representar el volumen de un prisma en términos de sus aristas.
  \item Realiza la representación gráfica del desarrollo plano de un prisma.
  \item Estima, calcula y compara volúmenes a partir de las relaciones entre las aristas de un prisma o de otros sólidos.
  \item Interpreta las expresiones algebraicas que representan el volumen y el área cuando sus dimensiones varían.
\end{itemize}

\ejemplo

En la figura se presentan cinco cajas y en la tabla se especifican las dimensiones de cada una de ellas. Completa e interpreta la tabla a partir de las dimensiones de cada caja.

\fig{g08_p61_1}{Hoja con una secuencia de figuras hechas con cubos y una tabla para registrar el patrón.}

Encuentra las razones aritméticas entre los diferentes volúmenes de las cajas y la expresión general para el volumen y el área exterior total de cada una de ellas.


% PDF p. 61
\dba{5}{Utiliza y explica diferentes estrategias para encontrar el volumen de objetos regulares e irregulares en la solución de problemas en las matemáticas y en otras ciencias.}

\evidencias

\begin{itemize}
  \item Estima medidas de volumen con unidades estandarizadas y no estandarizadas.
  \item Utiliza la relación de las unidades de capacidad con las unidades de volumen (litros, dm 3 , etc) en la solución de un problema.
  \item Identifica la posibilidad del error en la medición del volumen haciendo aproximaciones pertinentes al respecto.
  \item Explora y crea estrategias para calcular el volumen de cuerpos regulares e irregulares.
\end{itemize}

\ejemplo

En un recipiente cilíndrico totalmente lleno de agua, se sumerge por completo un objeto de forma irregular, el agua desalojada se recoge en un recipiente que se ha colocado previamente como lo muestra la figura.

\fig{g08_p61_2}{Persona llenando un recipiente con líquido (situación de variación).}

\fig{g08_p61_3}{Persona llenando un recipiente de otra forma (comparar la variación de la altura).}

Compara el volumen calculado con el volumen de la cantidad de agua derramada. Describe los procedimientos utilizados y explica los resultados y sus respectivos procedimientos. Asocia la forma del objeto irregular formada por una composición de figuras regulares, utiliza estas figuras para calcular el volumen del objeto irregular con una aproximación razonable.


% PDF p. 62
\dba{6}{Identifica relaciones de congruencia y semejanza entre las formas geométricas que configuran el diseño de un objeto.}

\evidencias

\begin{itemize}
  \item Utiliza criterios para argumentar la congruencia de dos triángulos.
  \item Discrimina casos de semejanza de triángulos en situaciones diversas.
  \item Resuelve problemas que implican aplicación de los criterios de semejanza.
  \item Compara figuras y argumenta la posibilidad de ser congruente o semejantes entre sí.
\end{itemize}

\ejemplo

Las grandes empresas invierten en el diseño de la imagen corporativa que los representa. Por ejemplo, en sus logotipos o iconos que los diferencian en el mercado. Las empresas buscan que los símbolos además de sencillos sean inconfundibles, para que las personas siempre los distingan entre las demás marcas. En muchos de los logotipos de las grandes marcas priman las regularidades geométricas como se muestra a continuación:

\fig{g08_p62_1}{Figura con simetría rotacional (hexágono estilizado).}

\fig{g08_p62_2}{Logotipo con simetría (arcos entrelazados).}

\fig{g08_p62_3}{Figura de tres pétalos con simetría rotacional de orden 3.}

Identifica las figuras congruentes que hay en cada uno de los logotipos. Argumenta las congruencias encontradas en cada logotipo.


% PDF p. 62
\dba{7}{Identifica regularidades y argumenta propiedades de figuras geométricas a partir de teoremas y las aplica en situaciones reales.}

\evidencias

\begin{itemize}
  \item Describe teoremas y argumenta su validez a través de diferentes recursos (Software, tangram, papel, entre otros).
  \item Argumenta la relación pitagórica por medio de construcción al utilizar material concreto.
  \item Reconoce relaciones geométricas al utilizar el teorema de Pitágoras y Thales, entre otros.
  \item Aplica el teorema de Pitágoras para calcular la medida de cualquier lado de un triángulo rectángulo.
  \item Resuelve problemas utilizando teoremas básicos.
\end{itemize}

\ejemplo

A partir de los rompecabezas que se muestran en la imagen, explica regularidades y propiedades que se presentan al variar los elementos de la construcción geométrica del teorema de Pitágoras hasta llegar a su generalización por medio de diferentes expresiones (numéricas, geométricas y algebraicas). Con el apoyo de Software verifica la generalidad del teorema de Pitágoras con triángulos, cuadrados y formas de diferentes tamaños.

\fig{g08_p62_4}{Cuadrado sobre cuadrícula con una diagonal trazada.}

\fig{g08_p62_5}{Cuadrado inclinado inscrito en otro cuadrado sobre cuadrícula (áreas y teorema de Pitágoras).}


% PDF p. 63
\dba{8}{Identifica y analiza relaciones entre propiedades de las gráficas y propiedades de expresiones algebraicas y relaciona la variación y covariación con los comportamientos gráficos, numéricos y características de las expresiones algebraicas en situaciones de modelación.}

\evidencias

\begin{itemize}
  \item Opera con formas simbólicas y las interpreta.
  \item Relaciona un cambio en la variable independiente con el cambio correspondiente en la variable dependiente.
  \item Encuentra valores desconocidos en ecuaciones algebraicas.
  \item Reconoce y representa relaciones numéricas mediante expresiones algebraicas y encuentra el conjunto de variación de una variable en función del contexto.
\end{itemize}

\ejemplo

Escribe una expresión que relacione el cambio que ocurre en el valor del volumen del cono circular recto cuando el radio cambia de r a r+ ∆r (∆r representa un incremento en el valor de r) y la altura permanece constante. Calcula el cambio en el volumen para algunos incrementos del radio. Representa por medio de una gráfica la relación entre volumen cuando la altura del cono y el radio de su base son iguales y la utiliza para averiguar el valor del radio para que el volumen sea igual a 20 u 3 (u es la unidad de medida). Encuentra los valores enteros del volumen a partir de los valores del radio.


% PDF p. 63
\dba{9}{Propone, compara y usa procedimientos inductivos y lenguaje algebraico para formular y poner a prueba conjeturas en diversas situaciones o contextos.}

\evidencias

\begin{itemize}
  \item Opera con formas simbólicas que representan números y encuentra valores desconocidos en ecuaciones numéricas.
  \item Reconoce patrones numéricos y los describe verbalmente.
  \item Representa relaciones numéricas mediante expresiones algebraicas y opera con y sobre variables.
  \item Describe diferentes usos del signo igual (equivalencia, igualdad condicionada) en las expresiones algebraicas.
  \item Utiliza las propiedades de los conjuntos numéricos para resolver ecuaciones.
\end{itemize}

\ejemplo

Encuentra valores para b, c, d, e, etc., que satisfagan las ecuaciones propuestas y argumenta cómo cambian las respuestas obtenidas si se cambia el valor de a por 6 o por 8.

a = 4 a + 2b = 10 a + 2b + 3c = 28 a + 2b + 3c + 4d = 68 a + 2b + 3c + 4d + 5e = 93 a +2b + 3c + 4d + 5e + 6f = 123 a + 2b + 3c + 4d + 5e + 6f + 7g = 200

Describe los procedimientos para obtener valores numéricos que satisfagan las ecuaciones segunda y tercera, si se desconoce el valor de a.


% PDF p. 64
\dba{10}{Propone relaciones o modelos funcionales entre variables e identifica y analiza propiedades de covariación entre variables, en contextos numéricos, geométricos y cotidianos y las representa mediante gráficas (cartesianas de puntos, continuas, formadas por segmentos, etc.).}

\evidencias

\begin{itemize}
  \item Toma decisiones informadas en exploraciones numéricas, algebraicas o gráficas de los modelos matemáticos usados.
  \item Relaciona características algebraicas de las funciones, sus gráficas y procesos de aproximación sucesiva.
\end{itemize}

\ejemplo

Una araña ubicada en una esquina quiere cazar a una mosca que está ubicada en la esquina inferior izquierda de una caja cúbica cuyo lado mide un metro. La araña usará un camino recto pasando por dos caras del cubo y atravesando una de sus aristas por un punto X como se muestra en la línea punteada de la figura. Determina la posición del punto X para que el camino seguido por la mosca sea el más corto. Encuentra el camino más corto que ha de seguir la araña para llegar hasta la mosca.

\fig{g08_p64_1}{Cubo con una araña en S y una mosca en F; el camino por las caras pasa por X (camino más corto).}

Utiliza el teorema de Pitágoras para obtener la distancia de la línea punteada. Explora numéricamente una solución y elabora un modelo algebraico de las longitudes de las rutas posibles.


% PDF p. 64
\dba{11}{Interpreta información presentada en tablas de frecuencia y gráficos cuyos datos están agrupados en intervalos y decide cuál es la medida de tendencia central que mejor representa el comportamiento de dicho conjunto.}

\evidencias

\begin{itemize}
  \item Interpreta los datos representados en diferentes tablas y gráficos.
  \item Usa estrategias gráficas o numéricas para encontrar las medidas de tendencia central de un conjunto de datos agrupados.
  \item Describe el comportamiento de los datos empleando las medidas de tendencia central y el rango.
  \item Reconoce cómo varían las medidas de tendencia central y el rango cuando varían los datos.
\end{itemize}

\ejemplo

Los estadísticos que hicieron un estudio sobre la producción de café por hectárea en 510 fincas cafeteras cometieron un error, no incorporaron los datos de 60 fincas de un municipio. Ellos afirman que, como en esa población la producción de café por hectárea se encuentra entre los límites menor y mayor de las ya estudiadas, en general los resultados no varían.

\fig{g08_p65_1}{Gráfica ilustrada de producción de café por hectárea.}

\fig{g08_p65_2}{Gráfica «Producción de café por hectárea»: arrobas por hectárea 101, 151, 201, 251, 301, 401 → producción 180, 274, 317, 379, 470, 510.}

Lee y compara la información presentada en cada gráfica. Encuentra las medidas de tendencia central adecuadas y analiza si hay cambios o no cuando se introduce la información faltante.


% PDF p. 65
\dba{12}{Hace predicciones sobre la posibilidad de ocurrencia de un evento compuesto e interpreta la predicción a partir del uso de propiedades básicas de la probabilidad.}

\evidencias

\begin{itemize}
  \item Identifica y enumera el espacio muestral de un experimento aleatorio.
  \item Identifica y enumera los resultados favorables de ocurrencia de un evento indicado.
  \item Asigna la probabilidad de la ocurrencia de un evento usando valores entre 0 y 1.
  \item Reconoce cuando dos eventos son o no mutuamente excluyentes y les asigna la probabilidad usando la regla de la adición.
\end{itemize}

\ejemplo

Se realiza un estudio con estudiantes de grado octavo para indagar por la cantidad de hermanos y sus edades. En la tabla se presentan la cantidad de estudiantes por cada número de hermanos y rango de edades.

\fig{g08_p65_3}{Tabla de frecuencias: número de hermanos según rangos de edad.}

Con base en la tabla de distribución de frecuencias determina:

\begin{itemize}
  \item La probabilidad de que un estudiante del curso tenga 1 o 2 hermanos.
  \item La probabilidad de que en la familia de un alumno del curso haya dos hijos y él sea el menor.
  \item Utiliza al menos dos procedimientos diferentes para calcular esta probabilidad y justifica la igualdad de los resultados.
\end{itemize}

